% Supplementary material: method-level detail for Section 4 of the main text.

\section{Compression Methods in Detail}
\label{supp:compression}

This section expands the summary of Section~\ref{sec:compression} of the main text.

\subsection{Post-Training Quantization}

Post-training quantization needs only a small calibration set (typically 128--512 samples) and no gradient computation, which is why it is the dominant paradigm. GPTQ \citep{frantar2023gptq} quantizes layer by layer, using an approximation of the inverse Hessian of the layer-wise reconstruction loss to choose rounding decisions. \gls{awq} \citep{lin2024awq} observes that about 1\% of weights, those on large-magnitude activation channels, account for a disproportionate share of the quantization error, and applies per-channel scaling to protect them; its authors report higher accuracy than GPTQ at matched bit-width. At 8 bits the obstacle is activations rather than weights: LLM.int8() \citep{dettmers2022llmint8} keeps outlier dimensions in 16-bit, and SmoothQuant \citep{xiao2023smoothquant} migrates the difficulty from activations to weights to enable W8A8 integer kernels; both report accuracy essentially unchanged from FP16. SqueezeLLM \citep{kim2024squeezellm} uses sensitivity-based non-uniform quantization with a dense-and-sparse decomposition. \gls{gguf} \citep{gerganov2023llamacpp} packages weights, vocabulary, and metadata in one file, with block-wise ``K-quant'' schemes from Q2\_K to Q8\_0; mixed schemes such as Q4\_K\_M keep the most sensitive tensors at higher precision.

Two independent benchmarks bear on small models specifically. ThinkSLM \citep{srivastava2025thinkslm} evaluates 72 small models, including quantized and pruned derivatives, on 17 reasoning benchmarks with three repeats, and finds that quantization preserves reasoning ability while pruning disrupts it. SLMQuant \citep{slmquant2025} evaluates quantization methods on two sub-1B models (SmolLM-135M and Qwen2.5-0.5B) on knowledge, commonsense, and language-modeling tasks, and finds that methods tuned on large models transfer imperfectly and that sensitivity depends on architecture: under W4A8 quantization the larger of its two models scored lower than the smaller, which the authors attribute to architectural differences rather than size. Neither study establishes that 4-bit loss grows monotonically as model size falls, and we make no such claim.

\subsection{Quantization-Aware and Low-Bit Training}

Quantization-aware training simulates quantization in the training loop so that weights become robust to it. It recovers part of the accuracy that post-training quantization loses at 2--3 bits, with the margin widening as bit-width falls, but needs training data and compute, which is why it appears mainly in vendor pipelines such as Apple's 2-bit on-device quantization \citep{gunter2025afm}. Its limit case is training in the low-bit format from the start: BitNet b1.58 \citep{ma2024bitnet158} restricts weights to $\{-1,0,1\}$, and the released 2B4T model was trained in that format on four trillion tokens.

\subsection{Pruning}

SparseGPT \citep{frantar2023sparsegpt} extends the GPTQ machinery to weight removal and prunes OPT-175B and BLOOM-176B to 50--60\% unstructured sparsity in one shot with a negligible increase in perplexity. Wanda \citep{sun2024wanda} removes the weights with the smallest product of magnitude and input-activation norm, with no weight update, and its authors report parity with SparseGPT on LLaMA and Llama~2; that parity does not hold below 4B \citep{kurz2025pruning_vs_quant}. Structured pruning removes heads, channels, or layers and yields models that run on dense kernels: LLM-Pruner \citep{ma2023llmpruner} removes coupled structures selected by gradient information and recovers quality with a short \gls{lora} tuning; ShortGPT \citep{men2024shortgpt} deletes the layers that change their input least, and we found no study establishing how far this tolerance holds below 7B (\evD). Sheared-LLaMA \citep{xia2024sheared} prunes Llama~2 7B to 1.3B and 2.7B and reports outperforming same-size open models at about 3\% of the compute of training them from scratch (\evC).
